% ---------------------------------------------------------------------
% Sec. VI-E of the manuscript, which is currently a heading with no body.
% Numbering continues from (105).
%
% This is the section that collects the two channel maxima.  The
% per-channel explanations were removed from VI-B and VI-C on purpose, so
% rho_phi,max and rho_GE,max are computed here once, and (105) is turned
% into a number.
%
% Every constant is taken from the manuscript as it stands:
%   l_p + p_off     0.140 + 0.020 = 0.160 m (roll), 0.110 + 0.020 = 0.130 (pitch)
%   |beta_M|        0.03446 (roll), 0.02573 (pitch)         -- table 3
%   W               30.08 N unloaded, 31.59 N ballasted
%   z_CoM           0.30 m, top of the Sec. VI-D box
%   J_CoM^CAD       0.0537 kg m^2
%   phi_max         10 deg
%   Mdot            0.10 to 1.20 N m/s                      -- table 2
%
% NOTE the earlier working used |beta_M| = 0.03446 on both axes.  With the
% pitch value from table 3 the pitch bound is 0.322 mm, not 0.331.
%
% Reproduced by analysis/vi_e_geometric_box.py.
% ---------------------------------------------------------------------

\documentclass[10pt]{article}
\usepackage[a4paper,margin=22mm]{geometry}
\usepackage{amsmath,amssymb,booktabs}
\setlength{\parskip}{0.35em}\setlength{\parindent}{0pt}
\newcommand{\Wz}{W\!z_{CoM}}
\newcommand{\lpo}{\bigl(l_p + p_{\rm off}\bigr)}
\title{\vspace{-16mm}\textbf{Sec.~VI-E: hand-calculable bounds}\vspace{-3mm}}
\date{\small numbering continues from (105)}
\begin{document}\maketitle

Eq.~(105) bounds the shifted critical moment by the two channel maxima,
and (93) supplies the constants that multiply them. What remains is to
evaluate $\rho_{\varphi,\max}$ and $\rho_{GE,\max}$ themselves. Both are
collected here rather than at the point where each channel was
introduced, so that the arithmetic appears once.

Nothing below requires a run. The weight comes from a scale, the pivot
arm and $z_{CoM}$ from the box of Sec.~VI-D, the inertia from CAD, the
tilt cap and the ramp rates from the excitation protocol of Sec.~V-B,
and the ground-effect slope from table~3. The bound is therefore a
statement made before the campaign, not a description of it.

\paragraph{The window, which only the coupling channel needs.}
The nominal excursion of (20) is
$\varphi_{\rm end} = \bigl(\dot M/\Wz C_2\bigr)\Lambda(x)$, and every
term of $\Lambda$ is positive, so $\Lambda(x)\ge x^3/6$. Imposing the
tilt cap $\varphi_{\rm end}\le\varphi_{\max}$ therefore inverts in
closed form, and with $J_P = \Wz/C_2^2$ the window length and the moment
increment across it follow without solving for $x$:
\begin{equation}
  \tau_{\rm end} \le
  \Bigl(\frac{6\varphi_{\max}J_P}{\dot M}\Bigr)^{1/3},
  \qquad
  \Delta M_{\rm win} = \dot M\tau_{\rm end}
  \le \bigl(6\varphi_{\max}J_P\dot M^{2}\bigr)^{1/3} .
  \tag{106}
\end{equation}
The inertia enters at its upper end, $J_P \le J_{CoM}^{\rm CAD} +
m\bigl(z_{CoM}^2 + \lpo^2\bigr)$, which is the unfavourable side here:
(82) discards that CAD term, but retaining it is what makes (106) a
bound rather than an estimate.

\paragraph{The gravity channel.}
By (85) the remainder is $\rho_\varphi = \tfrac12G''(\varphi^*)\varphi^2$
with $G''(\varphi^*) = W\lpo\cos\varphi^* - \Wz\sin\varphi^*$, which
decreases on $[0,\varphi_{\max}]$ and is therefore largest at
$\varphi^* = 0$. Hence
\begin{equation}
  \rho_{\varphi,\max} \;\le\; \tfrac12\,W\lpo\,\varphi_{\max}^{2},
  \tag{107}
\end{equation}
evaluated at the ballasted weight and at the widest arm the offset box
admits. It does not depend on the ramp rate: the tilt cap alone fixes
it.

\paragraph{The coupling channel.}
Term (iv) of (80) is $\rho_{GE}(\tau,\varphi) = \beta_M\dot M\tau\varphi$,
so its maximum over the window is the product of the moment increment
and the excursion, both already bounded:
\begin{equation}
  \rho_{GE,\max}
  = |\beta_M|\,\Delta M_{\rm win}\,\varphi_{\rm end}
  \;\le\; |\beta_M|\,\varphi_{\max}
          \bigl(6\varphi_{\max}J_P\dot M^{2}\bigr)^{1/3} .
  \tag{108}
\end{equation}
This is the only place the ramp rate enters, and it enters as
$\dot M^{2/3}$.

\paragraph{The bound, in six lines.}
Substituting (107) and (108) into (105) and dividing by the unloaded
weight gives the equivalent offset error. No hyperbolic function
appears, $C_2$ never enters, and $x$ is never solved for --- $\tfrac17$
and $\tfrac15$ in (93) are suprema over \emph{all} $x$:
\begin{equation}
  \boxed{\;
  \bigl|\Delta M_{{\rm crit},e}\bigr|
  \le \frac{\rho_{\varphi,\max}}{7} + \frac{\rho_{GE,\max}}{5},
  \qquad
  \bigl|\Delta\lambda_{\rm off}\bigr|
  \le \frac{1}{W}\Bigl[\frac{\rho_{\varphi,\max}}{7}
                     + \frac{\rho_{GE,\max}}{5}\Bigr].}
  \tag{109}
\end{equation}

Over the admissible box, with $W = 31.59$~N, $z_{CoM} = 0.30$~m,
$J_{CoM}^{\rm CAD} = 0.0537$~kg\,m$^2$, $\varphi_{\max} = 10^\circ$ and
the arms and slopes above:

\begin{center}\small
\begin{tabular}{@{}lrrrrrr@{}}
\toprule
axis & $\dot M$ & $J_P$ & $\Delta M_{\rm win}$
 & $\rho_{\varphi,\max}$ & $\rho_{GE,\max}$
 & $|\Delta\lambda_{\rm off}|$ \\
 & [N\,m/s] & [kg\,m$^2$] & [N\,m] & [mN\,m] & [mN\,m] & [mm] \\
\midrule
roll  & $0.10$ & $0.4260$ & $0.165$ & $76.98$ & $0.99$ & $0.372$ \\
      & $0.45$ & $0.4260$ & $0.449$ & $76.98$ & $2.70$ & $0.384$ \\
      & $1.20$ & $0.4260$ & $0.863$ & $76.98$ & $5.19$ & $0.400$ \\
\midrule
pitch & $0.10$ & $0.3979$ & $0.161$ & $62.55$ & $0.72$ & $0.302$ \\
      & $0.45$ & $0.3979$ & $0.439$ & $62.55$ & $1.97$ & $0.310$ \\
      & $1.20$ & $0.3979$ & $0.843$ & $62.55$ & $3.79$ & $0.322$ \\
\bottomrule
\end{tabular}
\end{center}

The worst admissible case is $0.400$~mm for roll and $0.322$~mm for
pitch, of which the gravity channel supplies $91$ and $92\%$. Two
readings follow. The inertia, the one input that is not directly
measured, reaches the result only through $\rho_{GE,\max}$ and only as a
cube root: halving $J_P$ moves the bound by $-1.8\%$ and doubling it by
$+2.2\%$, so a coarse CAD figure suffices. And precision is bought at
the tilt cap rather than at the ramp: (107) is quadratic in
$\varphi_{\max}$ while (108) is only $\dot M^{2/3}$, so a twelve-fold
reduction in ramp rate moves the bound by $7\%$ whereas halving the cap
quarters it. The $10^\circ$ used here was chosen for onset visibility
on the lightly loaded configurations.

\paragraph{The angular rate error.}
Eq.~(90) bounds the rate deviation itself, and unlike (109) it needs the
window. The same cubic inversion gives $x \le \bar x =
\bigl(6\varphi_{\max}\Wz C_2/\dot M\bigr)^{1/3}$, and $C_2$ has a
geometric ceiling in which the mass cancels,
\begin{equation}
  C_2^{2} = \frac{\Wz}{J_P}
  \;\le\; \frac{W z_{CoM}}{m\bigl(z_{CoM}^2 + \lpo^2\bigr)}
  = \frac{g\,z_{CoM}}{z_{CoM}^2 + \lpo^2} ,
  \tag{110}
\end{equation}
giving $5.05$~s$^{-1}$ for roll and $5.25$ for pitch. At the fastest
ramp $\bar x = 3.47$ and (90) returns $|e_\omega(\tau_{\rm end})| \le
0.103$~rad/s against a nominal $1.91$~rad/s, i.e.\ $5.4\%$; at the
slowest it returns $57\%$, because $\bar x$ enters through $\sinh$ and
the cubic shortcut is exponentially lossy there. That degradation does
not touch (109), which carries no $x$ at all --- which is why the
threshold bound, and not the rate bound, is what the offset budget
inherits.

\paragraph{What (109) covers.}
The displacement of the identified onset caused by the two forcing terms
the estimator cannot absorb, and nothing else. The baseline error of
(104), the departure of the applied moment from a linear ramp, and the
ground-contact geometry all lie outside it. Those are measured rather
than bounded a priori, and are reported with the realised budget in
Sec.~VIII-E, where the same runs give an actual propagated error some
fifty times inside (109).

\end{document}
